Los problemas de optimización de recursos no solo se presentan en la ingeniería, sino también en la vida cotidiana; un ejemplo representativo es la gestión del tiempo, un recurso escaso y de alto valor. Dado que los recursos son limitados, surge la necesidad de utilizarlos de manera eficiente, lo que da origen a los problemas de optimización. En términos generales, toda tarea u objetivo requiere de recursos para su realización, y la pregunta fundamental es cómo asignarlos de forma óptima.

Para ilustrar la problemática, consideremos el escenario de las plataformas de streaming: normalmente tienen un costo asociado para acceder a su catálogo, aunque muchas ofrecen periodos de prueba gratuitos. Si antes de que termine la suscripción se plantea seleccionar los animes que generen mayor satisfacción al verlos, considerando el tiempo y la energía disponibles —bajo la restricción de evaluar únicamente capítulos completos—, esta formulación constituye una instancia del problema de la mochila, un clásico problema de optimización.

A lo largo de este informe se analizarán las distintas heurísticas existentes para resolverlo, evidenciando que algunas entregan resultados consistentemente mejores que otras. En particular, se estudiará qué parámetros determinan realmente la calidad de las soluciones, así como los costos computacionales asociados y las técnicas de programación empleadas para la resolución del problema.
